\subsection{Proof Simplification: Task and Metrics} \label{sec:task-metrics}

\smallsec{Task Definition}
We formalize the proof simplification task as minimizing the complexity of a given proof. Specifically, for a valid formal statement $s$ with proof $p$, the goal is to produce an alternative proof $p^*$ of $s$ that minimizes a complexity measure $\mathcal{L}$:
\vspace{-6pt}
\begin{equation*}
p^* = \argmin_{x \text{ proves } s} \mathcal{L}(x)
\end{equation*}
%
Our method is agnostic to the choice of complexity measure $\mathcal{L}$, provided that it is deterministic and can be automatically computed from the proof. This flexibility encompasses the metrics used in prior work~\citep{ahuja2024improver}. Here, we adopt proof length as the measure of complexity, defined as the number of tokens produced by a Lean-specific tokenizer. Our proof length measure correlates with character count but does not penalize long identifier names, and it ignores comments and line breaks. We denote the length of a proof $x$ by $|x|$, i.e., $\mathcal{L}(x) = |x|$.

% The model is given as input a compiling (in Mathlib 4.19.0) Lean proof $x$. The objective is to output the shortest proof $y_x$ that correctly compiles and proves the same statement as $x$.

% We design a lexer specific to Lean in order to tokenize our proofs, and use the number of Lean tokens as our length metric. Compared to using character length, our length metric does not penalize long hypothesis or variable names. In addition, comments and newlines also do not affect the metric. We use $|x|$ to denote the length of a proof $x$.

\smallsec{Evaluation Metrics}
Given an original proof $p$ and $k$ candidate simplifications generated by the model, $p'_1, p'_2, \ldots, p'_k$, we define $l_i = \min(|p|, |p'_i|)$ if $p'_i$ is a valid proof and $l_i = |p|$ otherwise. (Intuitively, an invalid attempt reverts to the original proof length). We evaluate proof simplification using two metrics:

\begin{itemize}
\itemsep1pt
    \item $\min@k \triangleq \min_i \left\{ l_i \right\}$ denotes the minimum shortened proof length (lower is better).
    \item $\mathrm{red}@k \triangleq \max_i \left\{ \frac{|p|-l_i}{|p|} \right\} = 1 - \frac{\text{min@k}}{|p|}$ denotes the maximum relative proof length reduction from the original proof (higher is better).
\end{itemize}
%
Note that these metrics may not always be correlated: a method that only excels at shortening long proofs has a lower min@k and red@k than one that only excels at shortening short proofs. As with the pass@k metric \citep{chen2021evaluatinglargelanguagemodels}, we report our metrics via an unbiased estimator using $n > k$ samples. We average min@k and red@k across samples in a dataset to get overall length and reduction metrics.
